\section{Call to Action}

The Grounding Gap is not an inevitable property of all AI systems; it is a design choice embedded in the current paradigm of static, offline-trained models. Addressing it requires coordinated action from three communities.
% 落地鸿沟并非所有AI系统的必然属性；它是嵌入在当前静态、离线训练模型范式中的一种设计选择。解决它需要来自三个群体的协调行动。

\textbf{Researchers:} Shift evaluation focus from model-centric benchmarks to system-level diagnostics. Measure not just what the model knows (L1--L3), but how well the system observes L4 states and maintains belief over time. The three diagnostic questions in Section~5.3 provide a starting framework.
% \textbf{研究者：}将评估焦点从以模型为中心的基准转向系统级诊断。不仅衡量模型知道什么（L1--L3），还要衡量系统观测L4状态并随时间维护信念的能力。第5.3节中的三个诊断性问题提供了一个起始框架。

\textbf{System architects:} Build Harnesses as first-class components, not afterthoughts. \(\Omega\) (observation) and \(\Phi\) (belief update) should be designed with the same rigor as the LLM itself. Prioritize observation fidelity---sensors, APIs, connectors---over model scale once the prior is sufficiently rich.
% \textbf{系统架构师：}将框架作为一等组件来构建，而非事后的补丁。\(\Omega\)（观测）和\(\Phi\)（信念更新）应以与LLM本身相同的严谨程度来设计。一旦先验足够丰富，应将观测保真度——传感器、API、连接器——置于模型规模之上。

\textbf{Community:} Develop shared benchmarks and protocols for evaluating grounded agents. Current benchmarks test static reasoning; we need benchmarks that require real-time observation, belief maintenance, and recovery from unexpected changes. Such benchmarks would make the Grounding Gap measurable and progress accountable.
% \textbf{社区：}开发用于评估落地智能体的共享基准和协议。当前基准测试静态推理；我们需要要求实时观测、信念维护和从意外变化中恢复的基准。这样的基准将使落地鸿沟可衡量，使进展可问责。

The shift from larger static artifacts to dynamic, situated systems is not a technological inevitability---it is a choice. The sooner we make it, the sooner agentic AI will move from impressive demos to reliable action in the real world.
% 从更大的静态制品向动态、情境化系统的转变并非技术的必然——它是一种选择。我们越早做出这一选择，行动AI就能越早从令人印象深刻的演示走向现实世界中的可靠行动。